% \begin{remark}[Comparison with Steiner's general case~\cite{EPRINT:Steiner24b}]
%     Recall also Eq. (17) from~\cite{EPRINT:Steiner24b} for the worst case complexity estimate:
%     \[
%         \mathcal{O}\Bigg( m\cdot d^3\cdot \binom{n+d-1}{d}^{\omega+2} \Bigg)
%     \]
%     which, considering $m=\binom{n+d-1}{d}$, and $\omega=2$ becomes:
%     \[
%         \mathcal{O}\Bigg( m\cdot d^3\cdot \binom{n+d-1}{d}^{\omega+2} \Bigg) = \mathcal{O}\Bigg( d^3\cdot \Big(\frac{n}{n+d}\Big)^{5} \cdot \binom{n+d}{d}^{5} \Bigg).
%     \]
% \end{remark}


%%%%%%%%%%%%%%%%%%%%%%%%%%%%%
%%%%%%%%%%%%%%%%%%%%%%%%%%%%%


% \subsection{Solving LWE with Binary Errors}
% We now consider the case of Learning With Errors with binary noise (\belwe).
% This variant, first introduced in~\cite{MicPei13}, restricts the error distribution to the binary field (and is regarded as being as hard as standard LWE under comparable parameters --- find reference).
% The \belwe\ problem is particularly relevant in practice, as it enables more efficient implementations of LWE-based cryptographic schemes without sacrificing security.

% Formally, \belwe\ is an LWE instance over $\mathbb{Z}_q$ in which the error vector is sampled from $\{0,1\} ^m$~\cite{MicPei13}.
% A key advantage of this setting is that it allows for the straightforward construction of noise-free non-linear equations.
% Indeed, let $\mathbf{e} = (e_1,\dots,e_m) \in \{0,1\} ^m$ and define the polynomial $f(x) = x(x-1)$.
% Then $f(e_i)=0$ for all $1 \le i \le m$.

% Let $(G, \mathbf{s} \times G + \mathbf{e}) = (G, \mathbf{c}) \in \mathbb{Z}_q^{n \times m} \times \mathbb{Z}_q^{m}$ be sampled according to $L_{\mathbf{s},\mathcal{U}(\mathbb{F}_2)}^{(n)}$.
% For each $1 \le i \le m$, we can write
% \[
% e_i = c_i - \sum_{j=1}^n s_j G_{j,i}.
% \]
% It follows that the secret $\mathbf{s}\in\mathbb{Z}_q^n$ is a solution to the following algebraic system:
% \[
% f_i = f \bigg(c_i - \sum_{j=1}^n s_j G_{j,i} \bigg).
% \]

% The algorithm described in the previous section can be applied to this setting without modification.
% Using the established notation, the polynomial $f$ has degree $d = 2$.
% Consequently, in order to apply \Cref{alg:alg}, we require
% \[
% \binom{n+d-1}{d} = \binom{n+1}{2} = \frac{n(n+1)}{2} = \mathcal{O}(n^2)
% \]
% samples.
% Hence, the number of required samples grows polynomially in $n$.

% For $d = 2$, the condition required to apply the $S$-polynomial criterion is trivially satisfied.
% In particular, the relevant bound is satisfied (see \emph{[reference to the bound/relation to be satisfied]}).
% Moreover, only
% \[
% \binom{n+d-2}{d-1} = \binom{n}{1} = n
% \]
% $S$-polynomials need to be computed.

% The overall complexity of recovering the secret $\mathbf{s}$ is therefore
% \[
% \mathcal{O}\Bigg( \frac{1}{3}\binom{n+d}{d}^3 \Bigg) 
% = \mathcal{O}\Bigg(\frac{1}{3}\binom{n+2}{2}^3 \Bigg) 
% = \mathcal{O}\Bigg(\frac{1}{3} {\bigg(\frac{(n+2)(n+1)}{2} \bigg)}^3 \Bigg) 
% = \mathcal{O}\bigg(\frac{1}{24}\cdot n^6 \bigg)
% \]
% field operations.


%%%%%%%%%%%%%%%%%%%%%%%%%%%%%
%%%%%%%%%%%%%%%%%%%%%%%%%%%%%

\subsection{Solving LWE with Binary Secret}\label{sec:binarysecret}
In this \bslwe, the secret $\mathbf{s} = (s_1,\dots,s_n)$ satisfies $s_i \in \{0,1\}$ satisfies $x_i^2 - x = 0$, for $1 \le i \le n$, while the error distribution is left unchanged.
%Let $F = \{ x_1^2 - x_1,\dots,x_n^2 - x_n \}$ and let $F_{\mathrm{LWE}}$ denote the polynomial system associated with the LWE samples, where the corresponding error polynomial has degree $d$.

Thus, without loss of generality, we include the set $\{x_i^2 - x_i = 0: 1 \le i \le n\}$ of polynomials, in the polynomial system and reduce all polynomials modulo the ideal generated by the set. Consequently, each polynomial will contain monomials of the form $\mathbf{x}^\alpha = x_1^{\alpha_1}\cdots x_n^{\alpha_n}$, where $\alpha_i \in \{0,1\}.$
For polynomials of degree at most $d$, the number of such monomials is $\sum_{i=0}^d \binom{n}{i}$,
where the number of monomials of degree exactly $i$ is $\binom{n}{i}$.
% In particular, the proofs of \Cref{thm:result,thm:iterationCost} rely on the number of monomials of a given degree, which is modified in the present case.
The algorithm remains unchanged, and the complexity analysis carries over by replacing the new monomial count.

The number of samples required is $\binom{n}{d}$. Under this modification, the complexities are
\[
\mathscr{C}_{\texttt{Spoly}} = \mathcal{O} \left(\binom{n}{d}^3\right), \mathscr{C}_{\texttt{DiagA}} = \mathcal{O} \left(\binom{n}{d}^3\right),  \mathscr{C}_{\texttt{DiagB}}=\mathcal{O} \left(\binom{n}{d}^3\right)
\]

% \mathcal{O} \left(\frac{d}{2(n+d-1)} \binom{n}{d}^3\right)

% of  becomes 
% $\mathcal{O} \left(2\binom{n}{d}^3\right)$, the complexity of \texttt{DiagA} is $\mathcal{O} \left(\frac{1}{3}\binom{n}{d}^3\right)$ and the complexity of \texttt{DiagB} is $\mathcal{O} \left(\frac{d}{2(n+d-1)} \binom{n}{d}^3\right)$.
Combining these bounds as in \Cref{thm:result}, we obtain that the overall complexity of our algorithm in the \bslwe\ setting is $
\mathcal{O}\left(\dbinom{n}{d}^3\right)$ (field operations).



%%%%%%%%%%%%%%%%%%%%%%%%%%%%%
%%%%%%%%%%%%%%%%%%%%%%%%%%%%%


% \paragraph*{LWE with Hints}
% The Learning With Errors (LWE) problem with hints is a variant of the standard LWE problem in which the adversary is given additional auxiliary information, referred to as \emph{hints}, about the secret or the error. 
% These hints typically take the form of linear or nonlinear relations involving the secret vector and are intended to leak partial information while preserving security for appropriately chosen parameters.


% %%%%%%%%%%%%%%%%%%%%%%%%%%%%%
% %%%%%%%%%%%%%%%%%%%%%%%%%%%%%


% \paragraph*{Extension to Ring LWE}
% Describe how to extend the algorithm also in the case of RLWE, if possible.


